\section*{Hoofdstuk \ref{ch:b1:findim} --- Eindige dimensie}

\begin{solution}{exo:b1:findim:1}
\begin{enumerate}
  \item $z = -x - y$: $F = \{(x, y, -x-y)\} =
        \operatorname{Vect}\bigl((1,0,-1), (0,1,-1)\bigr)$; de twee
        vectoren zijn \omterm{def:b1:vspaces:free}{vrij} (\omterm{prop:b1:vspaces:coordinates}{coördinaten}): $\dim F = 2$.
  \item $G = \{(x, x, 2t, t)\} =
        \operatorname{Vect}\bigl((1,1,0,0), (0,0,2,1)\bigr)$: \omterm{def:b1:vspaces:free}{vrij},
        $\dim G = 2$.
  \item $P(1) = P'(1) = 0$ betekent $(X-1)^2 \mid P$
        (\cref{prop:b1:poly:multiplicity}): $P = (X-1)^2(aX + b)$.
        \omterm{def:b1:vspaces:free}{Basis} $\bigl((X-1)^2, X(X-1)^2\bigr)$, dimensie $2$.
\end{enumerate}
\end{solution}

\begin{solution}{exo:b1:findim:2}
$(3,5,7) = (1,1,1) + 2(1,2,3)$ en $(0,1,2) = (1,2,3) - (1,1,1)$:
beide zijn combinaties van de eerste twee, die \omterm{def:b1:vspaces:free}{vrij} zijn (niet
evenredig). \omterm{def:b1:findim:rank}{Rang} $2$; \omterm{def:b1:vspaces:free}{basis} van het \omterm{def:b1:vspaces:span}{opspansel}:
$\bigl((1,1,1), (1,2,3)\bigr)$.
\end{solution}

\begin{solution}{exo:b1:findim:3}
Probeer $e_1 = (1,0,0,0)$ en $e_3 = (0,0,1,0)$ toe te voegen. De
familie $\bigl((1,1,0,0), (0,0,1,1), e_1, e_3\bigr)$ is \omterm{def:b1:vspaces:free}{vrij}: een
nulcombinatie $\alpha(1,1,0,0) + \beta(0,0,1,1) + \gamma e_1 +
\delta e_3 = 0$ luidt $(\alpha + \gamma, \alpha, \beta + \delta,
\beta) = 0$, dus $\alpha = \beta = 0$, en daarna $\gamma = \delta =
0$. Vier \omterm{def:b1:vspaces:free}{vrije} vectoren in dimensie $4$: een \omterm{def:b1:vspaces:free}{basis}
(\cref{prop:b1:findim:twoofthree}).
\end{solution}

\begin{solution}{exo:b1:findim:4}
De graden $0, 1, 2, 3$ zijn paarsgewijs verschillend: \omterm{def:b1:vspaces:free}{vrij}
(\cref{prop:b1:vspaces:freecriteria}), vier vectoren in dimensie
$4$: een \omterm{def:b1:vspaces:free}{basis}. Voor $X^3$: ontwikkel naar beneden,
\[
X(X-1)(X-2) = X^3 - 3X^2 + 2X,
\qquad
X(X-1) = X^2 - X,
\]
dus $X^3 - X(X-1)(X-2) = 3X^2 - 2X$; en $3X^2 - 2X = 3(X^2 - X) + X
= 3\,X(X-1) + X$. Bijgevolg
\[
X^3 = X(X-1)(X-2) + 3\,X(X-1) + 1\cdot X + 0\cdot 1 :
\]
\omterm{prop:b1:vspaces:coordinates}{coördinaten} $(0, 1, 3, 1)$ op $\bigl(1, X, X(X-1),
X(X-1)(X-2)\bigr)$. (Dat zijn vermomde stirlinggetallen.)
\end{solution}

\begin{solution}{exo:b1:findim:5}
Grassmann: $\dim(F \cap G) = 4 + 5 - \dim(F + G)$, en $F + G$ is een
\omterm{def:b1:vspaces:subspace}{deelruimte} van $E$ die $G$ bevat: $5 \leq \dim(F+G) \leq 7$.
Bijgevolg $\dim(F \cap G) \in \{2, 3, 4\}$. Realisaties in $\R^7$
met canonieke \omterm{def:b1:vspaces:free}{basis} $(e_1, \dots, e_7)$, waarbij $G =
\operatorname{Vect}(e_1, \dots, e_5)$:
\begin{itemize}
  \item $F = \operatorname{Vect}(e_1, e_2, e_6, e_7)$: $F + G =
        \R^7$, doorsnede $\operatorname{Vect}(e_1, e_2)$, dimensie
        $2$;
  \item $F = \operatorname{Vect}(e_1, e_2, e_3, e_6)$: doorsnede van
        dimensie $3$;
  \item $F = \operatorname{Vect}(e_1, e_2, e_3, e_4) \subseteq G$:
        dimensie $4$.
\end{itemize}
\end{solution}

\begin{solution}{exo:b1:findim:6}
$H$ is een \omterm{def:b1:vspaces:subspace}{deelruimte} (\cref{exo:b1:vspaces:1}~(4), woordelijk).
Volgens de factorstelling (\cref{thm:b1:poly:factor}) geldt $P \in H
\iff P = (X-1)Q$ met $\deg Q \leq n - 1$: de \omterm{def:b1:logic:map}{afbeelding} $Q \mapsto
(X-1)Q$ is een lineaire bijectie van $\R_{n-1}[X]$ op $H$, dus is
een \omterm{def:b1:vspaces:free}{basis} van $H$
\[
\bigl((X-1),\ (X-1)X,\ (X-1)X^2,\ \dots,\ (X-1)X^{n-1}\bigr),
\qquad \dim H = n .
\]
Een \omterm{def:b1:vspaces:sum}{complementaire} rechte: $\operatorname{Vect}(1)$ (de constanten).
Inderdaad is $H \cap \operatorname{Vect}(1) = \{0\}$ (een constante
ongelijk aan nul wordt niet nul in $1$) en tellen de dimensies op
tot $n + 1$: volgens Grassmann is $H \oplus \operatorname{Vect}(1) =
\R_n[X]$.
\end{solution}

\begin{solution}{exo:b1:findim:7}
Verwijderen: laat men $u_p$ weg, dan kan het \omterm{def:b1:vspaces:span}{opspansel} alleen
krimpen; en het krimpt met hoogstens één dimensie, want $u_p$ weer
toevoegen aan een \omterm{def:b1:vspaces:free}{basis} van het kleinere \omterm{def:b1:vspaces:span}{opspansel} geeft een
voortbrengende familie van het grotere met hoogstens één extra
vector. Symmetrisch bij het toevoegen van een vector $v$: het
nieuwe \omterm{def:b1:vspaces:span}{opspansel} bevat het oude met hoogstens één extra
voortbrenger, dus is zijn dimensie $r$ (lag $v$ al in het
\omterm{def:b1:vspaces:span}{opspansel}) of $r + 1$ (anders, want volgens
\cref{prop:b1:vspaces:freecriteria}~(2) laat een \omterm{def:b1:vspaces:free}{basis} zich
uitbreiden). Beide \omterm{def:b1:logic:statement}{uitspraken} samen geven het besluit dat de \omterm{def:b1:findim:rank}{rang}
``$1$-lipschitziaans'' is.
\end{solution}

\begin{solution}{exo:b1:findim:8}
Langs een strikt stijgende keten stijgen de dimensies strikt ($F_i
\subseteq F_{i+1}$, $F_i \neq F_{i+1}$ en
\cref{thm:b1:findim:subspaces}: gelijkheid van de dimensies zou
gelijkheid van de ruimten afdwingen). Dus is $\dim F_1 < \dim F_2 <
\dots < \dim F_k$ een strikt stijgende rij gehele getallen in
$\intint{0}{n}$: hoogstens $n + 1$ waarden, dus $k \leq n + 1$. Een
maximale keten in $\R^n$:
\[
\{0\} \subsetneq \operatorname{Vect}(e_1) \subsetneq
\operatorname{Vect}(e_1, e_2) \subsetneq \dots \subsetneq \R^n ,
\]
van lengte precies $n + 1$.
\end{solution}

\begin{solution}{exo:b1:findim:9}
$F + G$ bevat $F$ strikt (want $G \not\subseteq F$), dus $\dim(F +
G) = n$ en geeft Grassmann $\dim(F \cap G) = (n-1) + (n-1) - n = n -
2$.

Algemene bewering, met inductie op $k$: de doorsnede $I_k$ van $k$
hypervlakken heeft $\dim I_k \geq n - k$. Waar voor $k = 1$. Stap:
$I_{k+1} = I_k \cap H_{k+1}$, en Grassmann binnen $E$ geeft
\[
\dim(I_k \cap H_{k+1}) = \dim I_k + (n - 1) - \dim(I_k + H_{k+1})
\geq \dim I_k + (n-1) - n \geq n - k - 1 . \qedhere
\]
\end{solution}

\begin{solution}{exo:b1:findim:10}
\begin{enumerate}
  \item De voorwaarde $u_{n+2} - 3u_{n+1} + 2u_n = 0$ is lineair en
        wordt door de nulrij vervuld: $E$ is een \omterm{def:b1:vspaces:subspace}{deelruimte}. Met
        inductie bepalen $u_0$ en $u_1$ elke $u_n$, en de
        afhankelijkheid is lineair (elke stap is een lineaire
        combinatie van de twee vorige waarden); omgekeerd komt elk
        paar $(a, b)$ van precies één rij van $E$ (definieer $u_n$
        met de recursie). Zoals in \cref{ex:b1:findim:dims} wordt
        $E$ \omterm{def:b1:logic:inj}{bijectief} en lineair geparametriseerd door $(u_0, u_1)
        \in \R^2$: $\dim E = 2$.
  \item Constanten: $3\cdot1 - 2\cdot1 = 1$. Meetkundig: $3\cdot
        2^{n+1} - 2\cdot 2^n = (6 - 2)2^n = 2^{n+2}$. Beide liggen
        in $E$. Vrijheid: $a\cdot 1 + b\cdot 2^n = 0$ voor alle $n$
        geeft, in $n = 0$ en $n = 1$: $a + b = 0$, $a + 2b = 0$, dus
        $a = b = 0$. Twee \omterm{def:b1:vspaces:free}{vrije} vectoren in dimensie $2$: een \omterm{def:b1:vspaces:free}{basis}
        (\cref{prop:b1:findim:twoofthree}).
  \item Los $a + b = 0$, $a + 2b = 1$ op: $b = 1$, $a = -1$, dus
        $u_n = 2^n - 1$ (de rij van Mersenne).
\end{enumerate}
\end{solution}

\begin{solution}{exo:b1:findim:11}
\begin{enumerate}
  \item Grassmann: $\dim(F \cap G) = \dim F + \dim G - \dim(F + G)
        \geq \dim F + \dim G - n > 0$, dus $F \cap G \neq \{0\}$.
        Met $n = 100$: $51 + 50 - 100 = 1 > 0$, de doorsnede bevat
        een rechte.
  \item Is $F \cap H = \{0\}$, dan is de som direct en geldt $\dim F
        + (n - 1) = \dim(F \oplus H) \leq n$, dus $\dim F \leq 1$.
        (Omgekeerd voldoet een rechte die niet in $H$ ligt hier wel
        aan: hypervlakken missen bijna niets.)
\end{enumerate}
\end{solution}

\begin{solution}{exo:b1:findim:12}
Neerwaartse inductie op $d = \dim F = \dim G$, van $d = n$ tot $d =
0$. Is $d = n$, dan is $F = G = E$ en voldoet $S = \{0\}$. Neem aan
dat de \omterm{def:b1:logic:statement}{uitspraak} geldt voor paren \omterm{def:b1:vspaces:subspace}{deelruimten} van dimensie $d + 1
\leq n$, en zij $\dim F = \dim G = d < n$.

Is $F = G$: neem voor $S$ een willekeurige \omterm{def:b1:vspaces:sum}{complementaire} van $F$
(\cref{thm:b1:findim:subspaces}). Is $F \neq G$: beide zijn echt,
dus bestaat er volgens \cref{exo:b1:vspaces:12} een $x \notin F \cup
G$. De sommen $F' = F \oplus \operatorname{Vect}(x)$ en $G' = G
\oplus \operatorname{Vect}(x)$ zijn direct ($x \notin F$, $x \notin
G$) en hebben dimensie $d + 1$; volgens de inductiehypothese laten
zij een gemeenschappelijke \omterm{def:b1:vspaces:sum}{complementaire} $S'$ toe: $E = F' \oplus
S' = G' \oplus S'$. Stel $S = \operatorname{Vect}(x) \oplus S'$ ---
direct, want $S' \cap \operatorname{Vect}(x) \subseteq S' \cap F' =
\{0\}$.

Dan is $F + S = F + \operatorname{Vect}(x) + S' = F' + S' = E$, en
$F \cap S = \{0\}$: is $f = \lambda x + s'$ met $f \in F$ en $s' \in
S'$, dan is $s' = f - \lambda x \in F' \cap S' = \{0\}$, dus $f =
\lambda x$, wat $\lambda = 0$ afdwingt ($x \notin F$) en $f = 0$.
Bijgevolg $E = F \oplus S$, en symmetrisch $E = G \oplus S$.
\end{solution}

\begin{solution}{pb:b1:findim:1}
\textbf{1.} $\Q$ bevat $0 \neq 1$, is stabiel onder optelling,
vermenigvuldiging en tegengestelden, en elk rationaal getal
ongelijk aan nul heeft een rationale inverse: een \omterm{def:b1:structures:field}{lichaam}
(\cref{def:b1:structures:field}). De definities van \omterm{def:b1:vspaces:span}{opspansel},
vrijheid en \omterm{def:b1:vspaces:free}{basis} en de bewijzen uit de hoofdstukken 18--19
gebruiken alleen de optelling van vectoren, de distributiviteit en
de rekenkunde van de scalairen in een \omterm{def:b1:structures:field}{lichaam} --- nooit een
absolute waarde, een ordening of een limiet. Dus is $\R$, met zijn
eigen optelling en de vermenigvuldiging $\Q \times \R \to \R$, een
$\Q$-vectorruimte, en gelden alle stellingen. Delen door een
scalair gebeurt één keer in het bewijs van
\cref{thm:b1:findim:exchange}: om ``naar $g_k$ op te lossen'' \omterm{def:b1:arith:divides}{deelt}
men door zijn coëfficiënt ongelijk aan nul.

\textbf{2.} Is $a + b\sqrt2 = 0$ met $a, b \in \Q$ niet beide nul,
dan zou $b \neq 0$ geven dat $\sqrt2 = -a/b \in \Q$, in tegenspraak
met de irrationaliteit van $\sqrt2$ (\cref{exo:b1:arith:7} met $p =
2$); dus $b = 0$ en daarna $a = 0$: \omterm{def:b1:vspaces:free}{vrij} over $\Q$. Over $\R$ is de
betrekking $\sqrt2\cdot 1 + (-1)\cdot\sqrt2 = 0$ niet-triviaal:
afhankelijk. Bijgevolg is $\dim_\Q \Q(\sqrt2) = 2$, met eenduidige
\omterm{prop:b1:vspaces:coordinates}{coördinaten} $(a, b)$.

\textbf{3.} $(a + b\sqrt2)(c + d\sqrt2) = (ac + 2bd) + (ad +
bc)\sqrt2 \in \Q(\sqrt2)$. Dan is
\[
\sigma(uv) = (ac + 2bd) - (ad + bc)\sqrt2
= (a - b\sqrt2)(c - d\sqrt2) = \sigma(u)\sigma(v),
\]
dus $N(uv) = uv\,\sigma(uv) = u\sigma(u)\,v\sigma(v) = N(u)N(v)$.
Was $N(u) = a^2 - 2b^2 = 0$ met $u \neq 0$, dan zou $b \neq 0$
geven dat $(a/b)^2 = 2$, een rationale vierkantswortel van $2$; dus
$b = 0$, en daarna $a = 0$ en $u = 0$: tegenspraak. Bijgevolg is
$N(u) \neq 0$ voor $u \neq 0$.

\textbf{4.} $u \cdot \dfrac{\sigma(u)}{N(u)} = \dfrac{N(u)}{N(u)} =
1$, en $\sigma(u)/N(u) \in \Q(\sqrt2)$: elk element ongelijk aan
nul is inverteerbaar binnen $\Q(\sqrt2)$, dat dus een deellichaam
van $\R$ is. Voorbeelden: $N(3 + 2\sqrt2) = 9 - 8 = 1$, dus
\[
\frac1{3 + 2\sqrt2} = 3 - 2\sqrt2 ;
\qquad
N(1 + \sqrt2) = -1, \quad
\frac1{1 + \sqrt2} = \sqrt2 - 1 .
\]

\textbf{5.} Is $\sqrt6 = p/q$, dan is $6q^2 = p^2$; links is de
exponent van $2$ gelijk aan $1 + 2v_2(q)$, oneven, en rechts aan
$2v_2(p)$, even: onmogelijk. Stel nu $\sqrt3 = a + b\sqrt2$ met $a,
b \in \Q$. Kwadrateren: $3 = a^2 + 2b^2 + 2ab\sqrt2$. Is $ab \neq
0$, dan is $\sqrt2 = (3 - a^2 - 2b^2)/(2ab) \in \Q$: onmogelijk. Is
$b = 0$, dan is $\sqrt3 = a \in \Q$, in tegenspraak met
\cref{exo:b1:arith:7} ($p = 3$). Is $a = 0$, dan is $\sqrt3 =
b\sqrt2$, en vermenigvuldiging met $\sqrt2$ geeft $\sqrt6 = 2b \in
\Q$: onmogelijk. Dus $\sqrt3 \notin \Q(\sqrt2)$.

\textbf{6.} De producten van de basisvectoren zijn
\[
\sqrt2\,\sqrt3 = \sqrt6,\quad
\sqrt2\,\sqrt6 = 2\sqrt3,\quad
\sqrt3\,\sqrt6 = 3\sqrt2,\quad
(\sqrt2)^2 = 2,\ (\sqrt3)^2 = 3,\ (\sqrt6)^2 = 6,
\]
alle in $M$. Een product van twee elementen van $M$ ontwikkelt door
bilineariteit tot rationale combinaties hiervan: $M$ is stabiel
onder vermenigvuldiging.

\textbf{7.} Zij $x + y\sqrt3 = 0$ met $x, y \in K = \Q(\sqrt2)$. Is
$y \neq 0$, dan is $\sqrt3 = -x/y \in K$ (vraag 4: $K$ is een
\omterm{def:b1:structures:field}{lichaam}), in tegenspraak met vraag 5. Dus $y = 0$, en daarna $x =
0$: $(1, \sqrt3)$ is \omterm{def:b1:vspaces:free}{vrij} over $K$. Zij brengt ook voort: $K +
K\sqrt3 = \{(a + b\sqrt2) + (c + d\sqrt2)\sqrt3\} =
\operatorname{Vect}_\Q (1, \sqrt2, \sqrt3, \sqrt6) = M$. Bijgevolg
is $\dim_K M = 2$.

\textbf{8.} Een betrekking $a + b\sqrt2 + c\sqrt3 + d\sqrt6 = 0$
hergroepeert zich als $(a + b\sqrt2) + (c + d\sqrt2)\sqrt3 = 0$ met
coëfficiënten in $K$; volgens vraag 7 worden beide nul, en volgens
vraag 2 is $a = b = 0$ en $c = d = 0$. De familie is dus \omterm{def:b1:vspaces:free}{vrij} en
$\dim_\Q M = 4$.

\textbf{9.} Elke $x \in M$ schrijft zich als $x = \sum_j y_j f_j$
met $y_j \in K$ (\omterm{def:b1:vspaces:free}{basis} van $M$ over $K$), en elke $y_j = \sum_i
\lambda_{ij} e_i$ met $\lambda_{ij} \in \Q$ (\omterm{def:b1:vspaces:free}{basis} van $K$ over
$\Q$); invullen geeft $x = \sum_{i,j} \lambda_{ij}\, e_i f_j$: de
producten brengen $M$ over $\Q$ voort.

\textbf{10.} Stel $\sum_{i,j} \lambda_{ij}\, e_i f_j = 0$ met
$\lambda_{ij} \in \Q$. Hergroepeer: $\sum_j \bigl(\sum_i
\lambda_{ij} e_i\bigr) f_j = 0$, en de binnenste sommen horen bij
$K$. De vrijheid van $(f_j)$ over $K$ geeft $\sum_i \lambda_{ij}
e_i = 0$ voor elke $j$; de vrijheid van $(e_i)$ over $\Q$ geeft dan
$\lambda_{ij} = 0$ voor alle $i, j$.

\textbf{11.} Volgens de vragen 9 en 10 is $(e_i f_j)_{i \leq m,\, j
\leq n}$ een \omterm{def:b1:vspaces:free}{basis} van $M$ over $\Q$ met $mn$ elementen:
\[
\dim_\Q M = m\,n = \dim_\Q K \cdot \dim_K M .
\]
Controle: $\dim_\Q K = 2$ en $\dim_K M = 2$ (de vragen 2 en 7)
geven $\dim_\Q M = 4$, en dat is vraag 8.

\textbf{12.} De vectoren $u a_1, \dots, u a_d$ horen bij $A$
(stabiliteit). Zij zijn \omterm{def:b1:vspaces:free}{vrij}: is $\sum_i \lambda_i\, u a_i = 0$,
dan is $u \sum_i \lambda_i a_i = 0$ in $\R$, en $u \neq 0$ dwingt
$\sum_i \lambda_i a_i = 0$ af, dus $\lambda_i = 0$ (de $a_i$ vormen
een \omterm{def:b1:vspaces:free}{basis}). Dus is $(u a_1, \dots, u a_d)$ een \omterm{def:b1:vspaces:free}{vrije familie} van
$d$ vectoren in $A$ met $\dim_\Q A = d$: een \omterm{def:b1:vspaces:free}{basis}
(\cref{prop:b1:findim:twoofthree}). In het bijzonder ontbindt $1
\in A$ als $1 = \sum_i \mu_i\, u a_i = u\,v$ met $v = \sum_i \mu_i
a_i \in A$: de inverse van $u$ ligt in $A$. Hiermee vindt men vraag
4 terug ($A = \Q(\sqrt2)$), en men bewijst in één klap dat $M$ een
deellichaam van $\R$ is (met vraag 6).

\textbf{13.} De $d + 1$ vectoren $1, x, x^2, \dots, x^{d}$ liggen
alle in $A$ (stabiliteit onder producten); een \omterm{def:b1:vspaces:free}{vrije familie} van
$A$ heeft hoogstens $d$ vectoren
(\cref{prop:b1:findim:twoofthree}), dus zijn zij afhankelijk: er
zijn rationale getallen $\lambda_0, \dots, \lambda_d$, niet alle
nul, met $\sum_k \lambda_k x^k = 0$. De \omterm{def:b1:poly:def}{veelterm} $P = \sum_k
\lambda_k X^k$ is ongelijk aan nul, heeft rationale coëfficiënten
en graad $\leq d$, en $P(x) = 0$.

\textbf{14.} $s^2 = 2 + 2\sqrt6 + 3 = 5 + 2\sqrt6$: \omterm{prop:b1:vspaces:coordinates}{coördinaten}
$(5, 0, 0, 2)$. Dan is
\[
s^3 = s\,s^2 = (\sqrt2 + \sqrt3)(5 + 2\sqrt6)
= 5\sqrt2 + 2\sqrt{12} + 5\sqrt3 + 2\sqrt{18}
= 11\sqrt2 + 9\sqrt3 ,
\]
\omterm{prop:b1:vspaces:coordinates}{coördinaten} $(0, 11, 9, 0)$ (met $\sqrt{12} = 2\sqrt3$ en
$\sqrt{18} = 3\sqrt2$). Ten slotte is $s^4 = (s^2)^2 = 25 +
20\sqrt6 + 24 = 49 + 20\sqrt6$: \omterm{prop:b1:vspaces:coordinates}{coördinaten} $(49, 0, 0, 20)$.

\textbf{15.} Stel $a + bs + cs^2 + ds^3 = 0$. Aflezen van de vier
\omterm{prop:b1:vspaces:coordinates}{coördinaten} op $(1, \sqrt2, \sqrt3, \sqrt6)$ geeft
\[
a + 5c = 0,\qquad b + 11d = 0,\qquad b + 9d = 0,\qquad 2c = 0 .
\]
Dus $c = 0$ en daarna $a = 0$; door de middelste vergelijkingen af
te trekken volgt $2d = 0$ en daarna $b = 0$: $(1, s, s^2, s^3)$ is
\omterm{def:b1:vspaces:free}{vrij}. Vier \omterm{def:b1:vspaces:free}{vrije} vectoren in $M$ van dimensie $4$: een \omterm{def:b1:vspaces:free}{basis}, dus
$\operatorname{Vect}_\Q(1, s, s^2, s^3) = M$. Uit vraag 14 volgt
$s^3 - 9s = 2\sqrt2$ en $11s - s^3 = 2\sqrt3$:
\[
\sqrt2 = \frac{s^3 - 9s}{2},
\qquad
\sqrt3 = \frac{11s - s^3}{2} .
\]

\textbf{16.} $s^4 - 10s^2 + 1 = (49 + 20\sqrt6) - 10(5 + 2\sqrt6) +
1 = 0$. Een \omterm{def:b1:poly:def}{monische veelterm} van graad $\leq 3$ die in $s$ nul
wordt zou een niet-triviale nulcombinatie van $(1, s, s^2, s^3)$
opleveren, in tegenspraak met vraag 15: $X^4 - 10X^2 + 1$ heeft de
kleinste graad. Haar wortels: $X^2 = 5 \pm 2\sqrt6 = (\sqrt3 \pm
\sqrt2)^2$, dus zijn de vier reële wortels $\pm(\sqrt3 + \sqrt2)$
en $\pm(\sqrt3 - \sqrt2)$, dat wil zeggen $\pm\sqrt2 \pm \sqrt3$.

\textbf{17.} Uit $s^4 - 10s^2 + 1 = 0$ volgt $s\,(10s - s^3) = 1$,
dus
\[
\frac1s = 10s - s^3 = 10(\sqrt2 + \sqrt3) - (11\sqrt2 + 9\sqrt3)
= \sqrt3 - \sqrt2 ,
\]
in overeenstemming met $(\sqrt3 + \sqrt2)(\sqrt3 - \sqrt2) = 1$. Is
$a + b\sqrt2 + c\sqrt3 + d\sqrt6 = r \in \Q$, dan is $(a - r) +
b\sqrt2 + c\sqrt3 + d\sqrt6 = 0$, en dwingt de vrijheid (vraag 8)
af dat $b = c = d = 0$ (en $a = r$). Voor $\sqrt2 + \sqrt3 +
\sqrt6$ zijn de \omterm{prop:b1:vspaces:coordinates}{coördinaten} $(0, 1, 1, 1)$ niet van die vorm:
irrationaal --- een \omterm{def:b1:logic:statement}{uitspraak} die sterker is dan
\cref{exo:b1:reals:7}, verkregen zonder enige kwadrateertruc.

\textbf{18.} Een rationale wortel $p/q$ (onvereenvoudigbaar) van
$X^3 - 2$ voldoet aan $p^3 = 2q^3$, en het criterium van
\cref{exo:b1:poly:5} geeft $p \mid 2$, $q \mid 1$: de kandidaten
$\pm1, \pm2$, waarvan de derde machten $\pm1, \pm8 \neq 2$ zijn.
Geen rationale wortel; in het bijzonder is $t = 2^{1/3} \notin \Q$,
dus is $(1, t)$ \omterm{def:b1:vspaces:free}{vrij} over $\Q$ (zoals in vraag 2).

\textbf{19.} Stel dat $(1, t, t^2)$ afhankelijk is: een zekere $P
\in \Q[X]$ ongelijk aan nul met $\deg P \leq 2$ heeft $P(t) = 0$;
kies zo'n $P$ van \emph{kleinste} graad $d \geq 1$. Volgens vraag
18 is $d \neq 1$, dus $d = 2$. Euclidische deling
(\cref{thm:b1:poly:division}): $X^3 - 2 = PQ + R$ met $Q, R \in
\Q[X]$ en $\deg R < 2$. Evaluatie in $t$: $0 = P(t)Q(t) + R(t) =
R(t)$, dus wordt $R$ nul in $t$; de minimaliteit van $d$ dwingt $R
= 0$ af. Dan is $X^3 - 2 = PQ$ met $\deg Q = 1$: de rationale
wortel van $Q$ is een rationale wortel van $X^3 - 2$, in
tegenspraak met vraag 18. Bijgevolg is $(1, t, t^2)$ \omterm{def:b1:vspaces:free}{vrij} en
$\dim_\Q A = 3$. Stabiliteit: $t^3 = 2$ herleidt elk product van
$1, t, t^2$ tot een combinatie ervan ($t\cdot t^2 = 2$, $t^2\cdot
t^2 = 2t$); $A$ bevat $1$: volgens vraag 12 is $A$ een deellichaam
van $\R$.

\textbf{20.} Stel $t \in M$. Dan is $t^2 \in M$ (stabiliteit), dus
$A \subseteq M$, en is $M$ een \omterm{def:b1:vspaces:def}{vectorruimte} over het \omterm{def:b1:structures:field}{lichaam} $A$:
de axioma's zijn die van de rekenkunde in $\R$, beperkt. Zij is
\omterm{def:b1:findim:def}{eindigdimensionaal} over $A$ (een eindige $\Q$-voortbrengende
familie brengt a fortiori voort over $A \supseteq \Q$). De
\omterm{pb:b1:findim:1}{torenwet} voor $\Q \subseteq A \subseteq M$ geeft
\[
4 = \dim_\Q M = \dim_\Q A \cdot \dim_A M = 3\,\dim_A M ,
\]
onmogelijk: $3$ \omterm{def:b1:arith:divides}{deelt} $4$ niet. Dus $2^{1/3} \notin M$: geen enkele
rationale combinatie van $1, \sqrt2, \sqrt3, \sqrt6$ is gelijk aan
$2^{1/3}$.

\textbf{21.} $\Q(\sqrt2) \subseteq M$, dus $t \notin \Q(\sqrt2)$.
Was $t + \sqrt2 = r \in \Q$, dan zou $t = r - \sqrt2 \in
\Q(\sqrt2)$: tegenspraak --- $2^{1/3} + \sqrt2$ is irrationaal. Was
$t = a + b\sqrt3$, dan zou $t \in \operatorname{Vect}_\Q(1,
\sqrt3) \subseteq M$: opnieuw een tegenspraak.

\textbf{22.} $B$ is een \omterm{def:b1:vspaces:def}{vectorruimte} over het \omterm{def:b1:structures:field}{lichaam} $A$
(beperking van de scalairen), en \omterm{def:b1:findim:def}{eindigdimensionaal} omdat $\dim_\Q
B$ eindig is. De \omterm{pb:b1:findim:1}{torenwet} voor $\Q \subseteq A \subseteq B$ geeft
$\dim_\Q B = \dim_\Q A \cdot \dim_A B$: de linkerfactor \omterm{def:b1:arith:divides}{deelt}.
Deellichamen van $M$ hebben dus $\Q$-dimensie $1$, $2$ of $4$:
dimensie $1$ is $\Q$ zelf, dimensie $4$ is $M$.

\textbf{23.} Zij $F \subseteq M$ een deellichaam met $\dim_\Q F =
2$, en $w \in F \setminus \Q$: $(1, w)$ is \omterm{def:b1:vspaces:free}{vrij}, en dus een \omterm{def:b1:vspaces:free}{basis}
van $F$. Volgens vraag 13 ($d = 2$) is $w^2 = \alpha + \beta w$
voor rationale $\alpha, \beta$. Met $v = w - \beta/2 \in F
\setminus \Q$:
\[
v^2 = w^2 - \beta w + \frac{\beta^2}4 = \alpha + \frac{\beta^2}4
\;\in\; \Q,
\]
en $F = \operatorname{Vect}(1, v)$. Schrijf nu $v = a + b\sqrt2 +
c\sqrt3 + d\sqrt6$ en werk uit:
\[
v^2 = (a^2 + 2b^2 + 3c^2 + 6d^2)
+ (2ab + 6cd)\sqrt2 + (2ac + 4bd)\sqrt3 + (2ad + 2bc)\sqrt6 .
\]
De drie irrationale \omterm{prop:b1:vspaces:coordinates}{coördinaten} worden nul: $ab + 3cd = ac + 2bd =
ad + bc = 0$. Is $a \neq 0$, dan is $b = -3cd/a$, en wordt de
tweede voorwaarde $c(a^2 - 6d^2) = 0$; omdat $a^2 = 6d^2$ met $d
\neq 0$ zou maken dat $\sqrt6 = \abs{a/d}$ rationaal is, geldt $c =
0$ of $d = 0$, en in beide gevallen dwingen de overige voorwaarden
$b = c = d = 0$ af, dus $v \in \Q$: uitgesloten. Dus $a = 0$, en de
voorwaarden luiden $3cd = 2bd = bc = 0$: hoogstens één van $b, c,
d$ is ongelijk aan nul. Bijgevolg is $v$ een rationaal veelvoud van
$\sqrt2$, $\sqrt3$ of $\sqrt6$, en
\[
F \in \bigl\{\Q(\sqrt2),\ \Q(\sqrt3),\ \Q(\sqrt6)\bigr\} ,
\]
en elk daarvan is inderdaad een $2$-dimensionaal deellichaam
(stabiliteit als in vraag 6, inversen volgens vraag 12): precies
drie kwadratische deellichamen.

\textbf{24.} $(1 + \sqrt2 + \sqrt3)(1 + \sqrt2 - \sqrt3) = (1 +
\sqrt2)^2 - 3 = 2\sqrt2$, dus
\[
\frac1{1 + \sqrt2 + \sqrt3}
= \frac{1 + \sqrt2 - \sqrt3}{2\sqrt2}
= \frac{\sqrt2 + 2 - \sqrt6}{4}
= \frac12 + \frac14\sqrt2 + 0\cdot\sqrt3 - \frac14\sqrt6 .
\]
(Controle: $(1 + \sqrt2 + \sqrt3)(2 + \sqrt2 - \sqrt6) = 4$ na
uitwerken.)

\textbf{25.} (i) De $\Q$-dimensie hangt aan elk deellichaam van
$\R$ een \emph{geheel getal}, en de \omterm{pb:b1:findim:1}{torenwet} laat die gehele
getallen langs insluitingen vermenigvuldigen: de dimensie gedraagt
zich als een rekenkundige invariant, en
deelbaarheidsvoorwaarden worden onmogelijkheidsbewijzen. (ii) De
\omterm{def:b1:findim:def}{eindige dimensie} dwingt elk element te voldoen aan een rationale
veeltermvergelijking ongelijk aan nul van graad hoogstens de
dimensie: eindigheid betekent algebraïciteit. (iii) De
twee-uit-drieregel maakte van ``vermenigvuldigen met $u$ beeldt een
\omterm{def:b1:vspaces:free}{basis} af op een \omterm{def:b1:vspaces:free}{vrije familie} van maximale omvang'' de
surjectiviteit, en bracht $u^{-1}$ zonder formule voort: de
inversen kwamen uit het tellen. (iv) De \omterm{pb:b1:findim:1}{torenwet} verbiedt een
$3$-dimensionaal \omterm{def:b1:structures:field}{lichaam} binnen een $4$-dimensionaal, en daarom is
$2^{1/3}$ niet vanuit $\sqrt2$ en $\sqrt3$ bereikbaar. De stelling
van deel III is de \emph{\omterm{pb:b1:findim:1}{torenwet} van Dedekind}, de openingszet van
de galoistheorie, die in het volume van bachelorjaar 3 wordt
ontwikkeld.
\end{solution}
